\section*{Глава 1. Анализ предметной области}
\addcontentsline{toc}{section}{Глава 1. Анализ предметной области}
\stepcounter{section}

\subsection{Обзор существующих программных продуктов по теме работы}

На современном рынке существует ряд решений для автоматизации управления
фитнес-центрами. Проведём сравнительный анализ наиболее распространённых
продуктов в таблице~\ref{tab:competitors}.

\begin{table}[ht!]
	\centering
	\caption{Сравнительный анализ существующих решений для фитнес-центров}
	\label{tab:competitors}
	\begin{tabular}{|l|p{0.3\textwidth}|l|}
		\hline
		\textbf{Название}  & \textbf{Ключевая \linebreak особенность}      & \textbf{Модель доступа} \\ \hline
		Mindbody           & Онлайн-запись и мобильное приложение          & SaaS (платная)          \\ \hline
		1С:Фитнес клуб    & Интеграция с бухгалтерией                     & Лицензия                \\ \hline
		Fitness Pro        & Управление абонементами                       & Подписка                \\ \hline
	\end{tabular}
\end{table}

Анализ показал, что большинство продуктов ориентировано на крупные сети
фитнес-клубов и обладает избыточным функционалом для небольших студий.
Это подтверждает необходимость разработки доступного веб-приложения,
адаптированного под потребности небольшого фитнес-центра.

\subsection{Анализ программных инструментов разработки веб-приложений}

Для реализации проекта рассматриваются современные инструменты fullstack-разработки.
Выбор обусловлен требованиями к производительности, масштабируемости и удобству
разработки:

\begin{itemize}
	\item \textbf{Frontend:} Выбран \textit{React} как наиболее распространённая
	      библиотека для построения интерактивных пользовательских интерфейсов
	      с богатой экосистемой готовых компонентов.
	\item \textbf{Backend:} Использование языка \textit{Go} совместно с фреймворком
	      \textit{Gin} обеспечивает высокую производительность, строгую типизацию
	      и удобное построение REST API.
	\item \textbf{ORM:} Библиотека \textit{GORM} упрощает работу с базой данных,
	      позволяя описывать модели данных в виде Go-структур и выполнять запросы
	      без написания сырого SQL.
	\item \textbf{База данных:} \textit{PostgreSQL} выбрана как надёжная реляционная
	      СУБД, поддерживающая сложные запросы и агрегацию данных.
	\item \textbf{Аутентификация:} Применение \textit{JWT-токенов} обеспечивает
	      безопасную авторизацию пользователей с разграничением ролей.
\end{itemize}

\subsection{Формулировка цели и задач работы}

На основе проведённого анализа существующих решений и инструментов разработки
была сформулирована цель исследования.

\textbf{Цель работы} --- проектирование и реализация веб-приложения для
фитнес-центра, обеспечивающего онлайн-запись клиентов на тренировки,
управление расписанием и личный кабинет пользователя с разграничением
прав доступа по ролям.

Для достижения поставленной цели необходимо решить следующие \textbf{задачи}:
\begin{enumerate}
	\item Провести анализ существующих программных продуктов и инструментов
	      разработки веб-приложений;
	\item Спроектировать модель данных и архитектуру приложения;
	\item Разработать систему аутентификации, авторизации и управления ролями
	      пользователей (клиент, тренер, администратор);
	\item Реализовать CRUD-интерфейс для управления тренировками, расписанием
	      и пользователями;
	\item Реализовать онлайн-запись на тренировки с проверкой доступных мест;
	\item Разработать личный кабинет пользователя с историей записей и
	      возможностью загрузки файлов;
	\item Реализовать агрегацию данных для формирования статистики посещаемости.
\end{enumerate}